\footnotesize%
\pstart%
\noindent%
\pagelabel{elevatio_intro}%
\indexN{\textls{Chauvin} (Chauvinus), Etienne 1640--1725}%
\indexN{\textls{Jablonski}, Johann Theodor 1654--1731}%
\indexN{\textls{Chuno} (Couneau, Cuneau), Johann Jacob Julius 1661--1715}%
Die im Folgenden edierten Texte N.\;\ref{60210} \textit{Annotatio ad hypothesin de ascensu vaporum} und N.\;\ref{37109+60209} \textit{De elevatione vaporum et de corporibus quae ob cavitatem inclusam in aere natare possint} sind ihrer Entstehung und ihrem Inhalt nach aufs Engste miteinander verbunden.
Ihren Ausgangspunkt stellt der Aufsatz \textit{Nova circa vapores hypohesis} dar, den Etienne Chauvin, Mitglied der 1700 gegründeten Brandenburgischen\indexO{Brandenburg} Sozietät der Wissenchaften zu Berlin,\indexO{Berlin} für die geplanten \textit{Collectanea}, später \textit{Miscellanea Berolinensia} -- die Zeitschrift der Sozietät -- beigesteuert hatte (vgl. \textit{MB} 1710, S.\;120--122).\cite{01415}\cite{01411}
Die \textit{Annotatio} N.\;\ref{60210} war -- oder zumindest begann als -- eine direkte Stellungnahme Leibnizens zu Chauvins Aufsatz, wie dies am Titel selbst erkennbar ist.
Aus ihr entstand in einem zweiten Schritt der Aufsatz N.\;\ref{37109+60209}, den Leibniz im ersten Band der Zeitschrift unmittelbar nach Chauvins Beitrag veröffentlichen ließ (vgl. \textit{MB} 1710, S. 123–128):\cite{01412}\cite{01411}
Zunächst verfasste er das stark bearbeitete Konzept \textit{L}, auf dem -- mittelbar oder unmittelbar -- die verschollene Abschrift beruhte, die für den Druck \textit{E\textsuperscript{1}} verwendet wurde (siehe die Synopsis der Überlieferung von N.\;\ref{37109+60209} weiter unten, S.\;\pageref{MB_1710_123-128+LH_37_03_120-122_uebrlief}).
Leibnizens Aufsatz war zwar, stillschweigend, eine Erwiderung auf Chauvins \textit{Nova circa vapores hypothesis}\cite{01415}, knüpfte aber auch an seine früheren eigenständigen Überlegungen zur Aerostatik an, wie die dreißig Jahre ältere Aufzeichnung \textit{Aero-nautica} (N.\;\ref{56215} im vorliegenden Band) zeigt.
Die Entstehung der Texte N.\;\ref{60210} und N.\;\ref{37109+60209} lässt sich anhand der zumeist brieflichen Quellen, die die Vorbeireitung des ersten Bandes der \textit{Miscellanea Berolinensia} dokumentieren, mit relativer Genauigkeit darlegen.
\pend%
\pstart%
Der Plan für den ersten Band der Zeitschrift wurde auf den \glqq Aktivitätskonferenzen\grqq\ der Sozietät geschmiedet, zu denen Leibniz im Winter 1706/1707 die Berliner\indexO{Berlin} Mitglieder zusammengerufen hatte (\textsc{Brather} 1993, S.\;279).\cite{01419} %
Auf der Konferenz vom 30.\;Dezember 1706 erklärte Chauvin seine Absicht, zur gemeinsamen Unternehmung mit einer verbesserten Fassung seiner alten Ansichten über das Schweben der Dünste beizutragen; das Sitzungsprotokoll rubriziert als Arbeitstitel \textit{Neue Hypothesis de elevatione vaporum} (\textsc{Brather} 1993, S.\;185\,f.).\cite{01419} %
Am 27.\;August 1707 teilte Johann Theodor Jablonski, Sekretär der Sozietät der Wissenschaften, dem Präses Leibniz mit, dass ein guter Teil der für die Zeitschrift vorgesehenen Beiträge bereits bei dem mit der redaktionellen Arbeit beauftragten Mitglied der Sozietät, Johann Jacob Julius Chuno, angekommen sei (\textit{LSB} I,\;27 N.\;252, S.\;428.3--7).\cite{02563}
Ob zu diesem Zeitpunkt auch Chauvin seinen Aufsatz \textit{Nova circa vapores hypothesis} gesendet hatte, ist nicht bekannt.
Getan hatte er das wohl aber, bevor Jablonski am 8.\;Oktober 1707 an Leibniz schrieb: \textit{Der H.\;Chovin dörfte wol nicht gerne sehen, wenn seine arbeit zurück gesezet bliebe} (\textit{LSB} I,\;27 N.\;299, S.\;498.15).\cite{02564}
Der Ausfsatz hatte offenbar Zweifel erweckt, vermutlich beim sachverständigen Chuno.%
\indexN{\textls{Chauvin} (Chauvinus), Etienne 1640--1725}%
\indexN{\textls{Jablonski}, Johann Theodor 1654--1731}%
\indexN{\textls{Chuno} (Couneau, Cuneau), Johann Jacob Julius 1661--1715}%
\pend%
\pstart%
Leibniz muss dann den Wunsch geäßert haben, selbst den Aufsatz zu prüfen, wenn ihm Chuno am 24.\;März 1708 folgendermaßen antwortete: \textit{On vous envoyera comme vous le desirez le discours de M\textsuperscript{r} Chauvin, pour faire les remarques que vous} \angles{\textit{aurez}} \textit{dessein d'y faire. Il est cepandant sur que si l'on n'insere pas ce discours dans les Miscellanea, que nous n'aurons plus rien de luy une autre fois} (LBr~185 Bl.\;77\,r\textsuperscript{o}).\cite{02565}
Der Aufsatz erreichte tatsächlich Leibniz als Beilage zum Brief des Sekretärs Jablonski vom 5.\;Mai 1708: \textit{Des Hrn.\;Chauvins dissertation, so mir der Hr.\;Chuno zu dem Ende zugestellet, kommt hiebei} (\textit{Berichte}, S.\;53).\cite{02566}
Somit kann Leibniz nicht vor Mitte Mai 1708 Gelegenheit gehabt haben, Chauvins Text zu lesen und die Abfassung der \textit{Annotatio ad hypothesin de ascensu vaporum} N.\;\ref{60210} zu beginnen.
Da das Konzept \textit{L} seines Aufsatzes N.\;\ref{37109+60209} \textit{De elevatione vaporum et de corporibus quae ob cavitatem inclusam in aere natare possint} textgenetisch von der \textit{Annotatio} N.\;\ref{60210} abhängt, ist Mitte Mai 1708 als Terminus post quem der Datierung beider Texte zu betrachten.
Das Konzept N.\;\ref{37109+60209}, \textit{L} ist sicherlich später als N.\;\ref{60210} entstanden; wie groß der zeitliche Abstand aber gewesen ist, das lässt sich nach heutigem Wissensstand nicht ermitteln.%
\indexN{\textls{Chauvin} (Chauvinus), Etienne 1640--1725}%
\indexN{\textls{Jablonski}, Johann Theodor 1654--1731}%
\indexN{\textls{Chuno} (Couneau, Cuneau), Johann Jacob Julius 1661--1715}%
\pend%
\pstart%
Chunos Feststellung gegenüber Leibniz am 6.\;April 1709 -- \textit{Nous attendrons donc ... L'hypothese des faveurs}\bracks{\textit{!}} \textit{de M\textsuperscript{r} Chauvin avec ce que vous y apporterez} (LBr~185 Bl.\;81\,v\textsuperscript{o})\cite{02567} -- zeigt, dass zu jenem Zeitpunkt der Aufsatz \textit{De elevatione vaporum} N.\;\ref{37109+60209} noch nicht Berlin\indexO{Berlin} erreicht hatte.
Dies könnte aber schon bald geschenen sein, etwa mit einer Paketsendung aus Hannover,\indexO{Hannover} deren Ankunft Jablonski am 27.\;April 1709 bestetigte (\textit{Berichte}, S.\;60).\cite{02568}
In seinem Schreiben an Leibniz vom 27.\;Juli 1709 stellte Chuno tatsächlich fest, dass das Setzen der die physikalischen und medizinischen Beiträge umfassenden zweiten Abteilung des \textit{Miscellanea}-Bandes bereits Fortschritte gemacht habe; allerdings fügt er hinzu: \textit{On ajoutera du reste vos pieces, cet-à-dire la description et dessein de la machine aritmetique et des Meteora là où il convient} (LBr~185 Bl.\;94\,r\textsuperscript{o}, 95\,v\textsuperscript{o}).\cite{02569}
Wenn Chuno hiermit auf Chauvins Aufsatz und Leibnizens Erwiderung anspielt, dann war Leibniz gegen Ende Juli 1709 seiner Sendung noch schuldig.
Diese muss aber spätestens wenige Tage danach -- etwa mit dem Paket, für das sich Jablonski bei Leibniz am 10.\;August 1709 bedankt (\textit{Berichte}, S.\;62)\cite{02570} -- angekommen sein; denn am 20.\;August 1709 teilte Chuno dem Präses der Sozietät mit, \textit{Nous sommes presentement aux Medica et physica dans l'impression des Miscellanea}, und erwähnt weder Chauvins noch Leibnizens Aufsätze "über den Aufstieg der Dünste" unter den Beiträgen, die für den Abschluss der physikalisch-medizinischen Abteilung dringend benötigt wurden (LBr~185 Bl.\;96\,r\textsuperscript{o}).\cite{02571}
Und am 10.\;September 1709 kündigt Chuno an: \textit{L'impression s'avance et nous avons presque fini les Physica} (LBr~185 Bl.\;98\,r\textsuperscript{o})\cite{02572} -- eine Aussage, die Jablonski vier Tage später bestätigen sollte (\textit{Berichte}, S.\;65).\cite{02573}
In der physikalisch-medizinischen Abteilung des gedruckten Bandes sind Chauvins und Leibnizens Aufsätze am Anfang anzutreffen, nämlich unter Nr.\;5 und Nr.\;6 (vgl. \textit{MB} 1710, Inhaltsverzeichnis, S.\;387).\cite{01411}
Daher müssen sie spätestens zu dem Zeitpunkt, als der Druck dieses frühen Abschnitts der Abteilung durchgeführt wurde, bei Chuno vorgelegen haben, d.\,h. spätestens Mitte August 1709 (wahrscheinlich waren sie aber bereits vor dem 27.\;April verfasst und versendet worden). %
Zu diesem Zeitpunkt muss Leibniz foglich sowohl das Konzept \textit{L} des Aufsatzes N.\;\ref{37109+60209} \textit{De elevatione vaporum} wie erst recht auch die \textit{Annotatio} N.\;\ref{60210} angefertigt haben.
Das Konzept \textit{L} diente als mittelbare oder unmittelbare Vorlage für die verschollene Abschrift (mit Leibnizens letzten Textänderungen), die nach Berlin verschickt und spätestens ab Mitte August 1709 für den Druck \textit{E\textsuperscript{1}} verwendet wurde.
\pend%
\pstart%
Der gedruckte erste Band der \textit{Miscellanea Berolinensia} wurde schließlich auf der Leipziger\indexO{Leipzig} Buchmesse im Frühjahr 1710 ausgestellt, wie Jablonski am 17.\;Mai desselben Jahres an Leibniz berichtet: \textit{Das Buch ist zwar in Leipzig schon heraus, hie aber wird es noch nicht ausgegeben, bis es am Hofe überreichet worden, wozu die Exemplaria bei dem Buchbinder in der Arbeit sind.} (\textit{Berichte}, S.\;75).\cite{02574}
Spätestens in der ersten Hälfte Mai 1710 muss demgemäß auch der Druck \textit{E\textsuperscript{1}} des Aufsatzes \textit{De elevatione vaporum} N.\;\ref{37109+60209} endgültig angefertigt worden sein.
\pend%
\pstart%
Aus dieser ausführlichen Darstellung ergibt sich die Zeitspanne von Mitte Mai 1708 bis spätestens Mitte August 1709 als Datierung der gesamten handschriftlichen Überlieferung der Texte N.\;\ref{60210} \textit{Annotatio ad hypothesin de ascensu vaporum} und N.\;\ref{37109+60209} \textit{De elevatione vaporum et de corporibus quae ob cavitatem inclusam in aere natare possint}, während der Druck N.\;\ref{37109+60209},\textit{E\textsuperscript{1}} bis zur ersten Hälfte Mai 1710 fertiggestellt wurde.%
\indexN{\textls{Chauvin} (Chauvinus), Etienne 1640--1725}%
\indexN{\textls{Jablonski}, Johann Theodor 1654--1731}%
\indexN{\textls{Chuno} (Couneau, Cuneau), Johann Jacob Julius 1661--1715}%
\pend%
\selectlanguage{latin}%
\frenchspacing%
\normalsize%
